\section{Integrity, Security, and Privacy}

\subsection{Integrity controls}

The release evidence uses hash-linked certificates, manifests, and optional ECDSA-P256-SHA256
signatures. The runtime image is immutable by digest, the intake binds the product commit and
configuration, and the release inventory binds artifact coordinates and member hashes. The paper
build refuses a mismatched product SHA, bundle hash, artifact digest, backend identity, branch,
quality disposition, robustness plan, or claim boundary.

The public companion is deterministic: members are sorted, timestamps derive from the source
commit epoch, paths are canonical, symbolic links and traversal paths are rejected, and
\texttt{MANIFEST.json} records SHA-256 and size for every member. Its verifier rechecks exact
membership and hashes before parsing evidence.

\subsection{Application and deployment boundary}

The product code forbids generic outbound HTTP clients under the application package, except for
explicitly gated metering behavior. This is an application invariant, not a firewall. Operators must
still enforce network egress policy, private subnets or equivalent isolation, TLS termination,
identity and access management, secret storage, patching, monitoring, backup, and retention.

Evidence can be encrypted at rest in the workflow and integrity checked through its sidecars.
Encryption keys, signing keys, runner credentials, and private evidence are absent from this paper
and companion. Public provenance includes only reviewed hashes and sanitized aggregate data.

\subsection{Privacy interpretation}

The native exact-output wall found zero exact matches to training rows over its declared comparison
columns in both 10,000-row branch outputs. This is a useful deterministic control against exact tuple
reproduction. It does not establish resistance to near-neighbor linkage, membership inference,
attribute inference, or reconstruction.

The quality certificate's privacy component is labelled a \emph{privacy quality heuristic}. No
formal differential-privacy mechanism or $\varepsilon$ budget is claimed. The release backend
explicitly declares \texttt{not\_differential\_privacy}. Any deployment handling personal data must
perform its own privacy impact assessment, data-minimization review, access-control design, and
retention decision.

\clearpage
\subsection{Threats considered and residual risks}

\begin{table}[htbp]
\centering
\footnotesize
\AccessibleTableHeaderRow
\begin{tabular}{@{}>{\RaggedRight\arraybackslash}p{0.24\linewidth}
                    >{\RaggedRight\arraybackslash}p{0.34\linewidth}
                    >{\RaggedRight\arraybackslash}p{0.32\linewidth}@{}}
\toprule
Threat & Control represented in this evidence & Residual risk \\
\midrule
Artifact substitution & Exact tag, commit, run, attempt, artifact ID, digest, size, member hash & Compromise outside reviewed trust and custody boundaries \\
Metric/report drift & Canonical metrics, PDF parity checks, companion projection & Unreviewed future renderer or schema changes \\
Unsafe archive content & Canonical paths, no duplicates, no links, exact membership & Vulnerabilities in external ZIP/PDF consumers \\
Private path or data disclosure & Sanitized robustness projection and content scan & Aggregate disclosure and inference require contextual review \\
Exact training-row reproduction & Final exact-output wall & Near duplicates and inference attacks not measured here \\
Application egress & Static no-generic-client gate & Deployment networking remains operator-controlled \\
Misreading missing metrics & Explicit unavailable/not-informative states & Human readers may still overgeneralize from available screens \\
\bottomrule
\end{tabular}
\caption{Selected threat controls and residual risks for the public technical artifact.}
\label{tab:threats}
\end{table}
